% Part: incompleteness
% Chapter: theories-computability
% Section: complete-decidable

\documentclass[../../../include/open-logic-section]{subfiles}

\begin{document}

\olfileid{inc}{tcp}{cdc}

\olsection{\printtoken{S}{axiomatizable} సంపూర్ణ సిద్ధాంతాలు నిర్ణయించదగినవి}

ప్రతి వాక్యం~$!A$కు $!A$ గానీ $\lnot !A$ గానీ నిరూపించదగినదైతే, ఆ
సిద్ధాంతాన్ని {\em సంపూర్ణం} అంటాం.

\begin{lem}
సిద్ధాంతం~$\Th{T}$ సంపూర్ణమూ \tetoken{స్వీకృతీకరించదగినదీ}{!!{axiomatizable}} అనుకుందాం. అప్పుడు
$\Th{T}$ నిర్ణయించదగినది.
\end{lem}

\begin{proof}
$\Th{T}$ సంపూర్ణమనీ, $A$ గణనీయ స్వీకృతాల సమితి అనీ అనుకుందాం.
$\Th{T}$ వైరుధ్యమైనదైతే అది స్పష్టంగా గణనీయం. (అల్గోరిథం: ``అవును అని
చెప్పు.'') కాబట్టి $\Th{T}$ అవైరుధ్యమైనది కూడా అని అనుకోవచ్చు.

ఒక వాక్యం~$!A$, $\Th{T}$లో ఉందో లేదో నిర్ణయించడానికి, $\Th{T}$ నుంచి
$!A$కు \tetoken{ఒక వ్యుత్పత్తి}{!!a{derivation}} కోసం, $\lnot !A$కు \tetoken{ఒక వ్యుత్పత్తి}{!!a{derivation}} కోసం ఏకకాలంలో
వెతకండి. $\Th{T}$ సంపూర్ణం కనుక వాటిలో ఒకటి తప్పక దొరుకుతుంది; $\Th{T}$
అవైరుధ్యమైనది కనుక, $\lnot !A$కు \tetoken{ఒక వ్యుత్పత్తి}{!!a{derivation}} దొరికితే $!A$కు
\tetoken{వ్యుత్పత్తి}{!!{derivation}} ఏదీ ఉండదు.

వేరే మాటల్లో, $\Th{T}$ \tetoken{గ.లె.}{!!{c.e.}} అని మనకు ఇప్పటికే తెలుసు. కాబట్టి మనం
ముందు నిరూపించిన ఒక సిద్ధాంతం వల్ల, $\Th{T}$ యొక్క పూరకం కూడా
\tetoken{గ.లె.}{!!{c.e.}} అని చూపితే చాలు. కానీ \tetoken{ఒక సూత్రం}{!!a{formula}}~$!A$, $\Th{\bar T}$లో ఉన్నప్పుడు,
అప్పుడు మాత్రమే $\lnot !A$, $\Th{T}$లో ఉంటుంది; అందువల్ల
$\Th{\bar T} \leq_m \Th{T}$.
\end{proof}

\end{document}
